\documentclass[10pt,a4paper]{article}
\usepackage{fontspec}
\usepackage{polyglossia}
\setmainlanguage{russian}
\setmainfont{Liberation Serif}
\usepackage[left=1.1cm,right=1.1cm,top=0.9cm,bottom=0.8cm]{geometry}
\usepackage{array}
\setlength{\parindent}{0pt}
\pagestyle{empty}
\renewcommand{\arraystretch}{1.25}

\begin{document}

\noindent{\Large\bfseries Правда, ложь и перебор \ · \ ответы}\hfill{\small 15 сентября. Детям не выдаётся}
\vspace{4pt}\hrule\vspace{8pt}

{\small Варианты A и B --- одна и та же задача со словарной заменой. Ответ по структуре
совпадает, записывается другими словами. Исключение --- задача 4: там замена слов ответа
не меняет, списать можно. Звёздочки стоят у трёх реально тяжёлых: 10$^*$, 11$^*$, 12$^*$.}

\vspace{6pt}

\begin{tabular}{@{}c|p{8.1cm}|p{8.1cm}@{}}
\textbf{№} & \hfil\textbf{вариант A}\hfil & \hfil\textbf{вариант B}\hfil\\ \hline
1 & \textbf{(а) 99 \quad (б) 9} & \textbf{(а) 59 \quad (б) 6}\\ \hline
2 & \textbf{Иван рыцарь, Пётр лжец} & \textbf{Илья рыцарь, Павел лжец}\\ \hline
3 & \textbf{12 конфет} \ ($\{12,13\}$, чётное) & \textbf{14 конфет} \ ($\{14,15\}$, чётное)\\ \hline
4 & \textbf{нет · да · нет} & \textbf{нет · да · нет} \ (тот же)\\ \hline
5 & \textbf{рыжий.} При этом в первом дупле орехов нет, а во втором есть & \textbf{полосатый.} При этом в первой норке желудей нет, а во второй есть\\ \hline
6 & Аня --- черепаха \newline Света --- кошка \newline \textbf{Маша --- собака} (она и сказала правду) & Оля --- ёжик \newline Рита --- хомяк \newline \textbf{Даша --- попугай} (она и сказала правду)\\ \hline
7 & \textbf{(а)} нет \ \textbf{(б)} самая хитрая --- \textbf{Инесса}, \newline порядок: Инесса, Алиса, Лариса & \textbf{(а)} нет \ \textbf{(б)} самая ленивая --- \textbf{Нюша}, \newline порядок: Нюша, Муся, Ляля\\ \hline
8 & \textbf{5 рыцарей} (и 5 лжецов) & \textbf{4 рыцаря} (и 4 лжеца)\\ \hline
9 & Алефтина --- черепаха --- парк \newline Бэлла --- свинка --- лес \newline Светлана --- собака --- стадион \newline Вероника --- кошка --- двор & Марина --- ёжик --- сквер \newline Полина --- хомяк --- поле \newline Тамара --- кролик --- каток \newline Юля --- попугай --- балкон\\ \hline
10$^*$ & \textbf{12 хулиганов} (и 13 отличников) & \textbf{10 хулиганов} (и 11 отличников)\\ \hline
11$^*$ & \textbf{оба гномы} & \textbf{оба тролли}\\ \hline
12$^*$ & \textbf{все десять заказали кофе} & \textbf{все двенадцать заказали кофе}\\
\end{tabular}

\vspace{9pt}\hrule\vspace{6pt}

\textbf{Задача 1.} Отрицание «самое меньшее 100» --- это весь набор $\{99, 98, 97, \ldots\}$,
спрашивается наибольшее из него. Ответ «меньше 100» верен по сути --- засчитывать.

\vspace{2pt}
\textbf{Задача 4.} Рыцарь сказал «все трёхголовые --- умные», но НЕ сказал «только трёхголовые».
Поэтому (1) и (3) неверны, а (2) верно: трёхголовое обязано быть умным, значит глупое
трёхголовым быть не может.

\vspace{2pt}
\textbf{Задача 8.} Если рыцарей $r$, то правду говорят ровно те, чей номер больше $r$;
их должно быть ровно $r$ штук, отсюда $r$ --- половина сидящих.

\vspace{2pt}
\textbf{Задача 5.} Если бы бельчонок был серым, обе фразы были бы ложны: из первой следовало бы,
что в другом дупле орехи ЕСТЬ, из второй --- что орехов нет нигде. Противоречие, значит он рыжий,
и тогда орехи только во втором дупле. Здесь важно, что отрицание «хотя бы в одном есть» --- это
«ни в одном нет», а не «ровно в одном нет».

\vspace{2pt}
\textbf{Задача 7.} Если бы самой хитрой была Алиса, она солгала бы, сказав «я хитрее Ларисы»~---
значит Лариса хитрее Алисы, и Алиса не самая. Противоречие, отсюда пункт (а). Дальше:
Лариса тоже не самая (иначе её фраза «Алиса не самая хитрая» ложна, то есть самая --- Алиса).
Остаётся Инесса; её фраза «Алиса хитрее меня» ложна, и всё сходится.

\vspace{2pt}
\textbf{Задача 10$^*$.} Хулиган, стоящий рядом с Владом, врёт --- значит между ним и Владом
хулиганов НЕ шесть. Разбирается счётом по обе стороны от 13-го места; ответ проверен полным
перебором всех расстановок на 25 местах.

\vspace{2pt}
\textbf{Задача 11$^*$ --- остроумный гроб, механика нестандартная.} Правдивость тут зависит
от ТЕМЫ высказывания, а не только от того, кто говорит. Если А эльф, он говорит про своё золото,
а не про гномов, значит правду --- тогда Б, сказавший «ты лжёшь», лжёт; но Б-гном лжёт только
про своё золото, а эльф --- только про гномов, и оба случая отпадают.
Значит А гном, он солгал про своё золото, а Б сказал правду --- и это возможно только если
Б гном (эльф, говоря про гнома А, обязан был бы солгать). Оба гномы. Перебрано машиной
по всем четырём сочетаниям.

\textbf{Задача 12$^*$ --- главный гроб листка.} Логик отвечает «нет», только если сам взял чай:
тогда он точно знает, что «всем кофе» неверно. Значит каждое «не знаю» означает, что сам
говорящий выбрал кофе, но про остальных не знает. Первые девять сказали «не знаю» --- все
девять взяли кофе. Десятый слышал их ответы, знает свой выбор, и раз он смог ответить «да»,
он тоже взял кофе. Ответ: все десять. Перебрано машиной по всем раскладам.

\vspace{3pt}
\textbf{Что объявить устно, в листке этого нет.} Рыцарь всегда говорит правду, лжец всегда
лжёт, отличить их с виду нельзя. Туристов в задачах этого листка нет.

\vspace{10pt}\hrule\vspace{5pt}
{\small Источники: Бураго, «Дневник математического кружка», занятия 2 и 3 --- задачи 1, 2, 3, 4, 5;
Знаменская, «100 задач»: 4 класс №1 (задача 6), 5 класс №83 (задача 8);
Торлопов, кружок МЦНМО 7--8 класс, 25.10.2019 --- задача 8;
Всеросс., 2015, ШЭ, 6.3 --- задача 7; Всеросс., 2023, ШЭ, 5.2 --- задача 10$^*$;
«Бельчонок», 2022, 5.2 --- задача 5; Математический праздник, 2010, 6.2 --- задача 11$^*$
(обе из подборки Ксении);
листок №24 «Логика» (179 школа, Дориченко и Шашков) --- задача 12$^*$.
Все ответы проверены программой 2026-09-14.}

\end{document}
